%!TEX root =  main.tex


\para{Speculative decoding (SD).}
We consider the standard SD setting, which involves a small draft model (the drafter) and a large target model (the verifier). 
At each SD step, the drafter first generates a sequence of draft tokens, after which the verifier evaluates them in parallel to decide whether to accept or reject them (i.e., verification), and then produces the final output. 
Standard SD does not allow parallelization between drafting and verification, since the verification inherently depends on the prior availability of the draft tokens. 

\para{Parallel speculative decoding (PSD).}
We consider a variant of PSD that systematically parallelizes the drafting and verification stages. By overlapping these stages, the draft model generates tokens for one batch of requests while the target model simultaneously verifies the draft tokens from another batch, thereby effectively hiding drafting latency and providing a significant speedup over standard SD.


\para{Notations.}
We now introduce the notation used throughout this paper. We consider an LLM inference service provider with a fixed capacity $C$, representing the maximum number of tokens its computing resources can process in parallel at each step. 
In practice, this capacity depends on the service provider's server configuration.
At each decoding step, the server processes a batch of $m$ requests $r_1, r_2, \cdots, r_m$.
Each request $r_i$ is associated with a partially generated sequence $S_{i,t_i} = (d_{i,1}, \ldots, d_{i, t_i})$, where $t_i$ represents the number of tokens generated so far for request $r_i$. 
We allow a variable draft window size for each request, i.e., $k_i$ for request $r_i$. 
At step $s$, the draft model $\cS$ drafts and selects a set $\cD_s \coloneqq \{(i, j) | i \in [m], j \in [t_i+k_i]\}$, such that the total number of selected tokens satisfies $|\cD_s| = \sum_{i=1}^{m} k_i = C$.


\para{Average throughput.} 
Let $\vone_{i,j}$ be an indicator variable that equals $1$ if the token $(i,j) \in D_s$ is verified, $0$ otherwise.
The average throughput is defined as the total number of verified tokens generated over $T$ steps, divided by the total time $\tau$ required to complete these steps:
$
    \textstyle \cG \coloneqq \tau^{-1}\sum_{s=1}^{T} {(\sum_{(i,j) \in \cD_s}\vone_{i,j}} + m ).
$


\para{End-to-end latency (E2EL).}
The end-to-end latency is defined as the average time from request arrival to completion for a set of requests under the SD inference system. 